% ===================================================================
%  TABLICA Nr 10 — Morze. — Port.
%  Meme regime que les precedents. Tableau a UNE SEULE COLONNE : les
%  cles vont t10-01-1 a t10-25-2, sans c1 ni c2, et l'apparat porte
%  le titre et le premier intertitre ensemble — il n'y a pas de
%  t10-tit-1.
%
%  UNE PHRASE ENTIERE EST EN GRAS AU 5, dans le fac-simile ido comme
%  dans le neerlandais : « Le spectacle de cette flotte, quand elle
%  manoeuvre pour lever l'ancre, est des plus interessants. » Ce
%  n'est pas un terme de planche, c'est le compositeur qui a laisse
%  courir son gras. On le reproduit, ponctuation comprise, parce que
%  les autres colonnes le reproduisent.
%
%  ET LE TABLEAU DISTINGUE DEUX SORTES DE DOCK : « flotacanta doki »
%  (27) sont les docks flottants ou l'on repare, « flot-doko » (53)
%  le bassin a flot ou les barques viennent s'abriter. Le polonais
%  les separe aussi — doki plywajace et dok postojowy — sinon deux
%  renvois de la planche tombent sur le meme mot.
%
%  « rempari » (15) sont les remparts, non les ramparts d'un navire ;
%  « kanaleto » le chenal balise, non un petit canal ; et
%  « brizanti » les brisants, distincts de « ondego » (92) qui est la
%  lame. La planche montre les trois.
% ===================================================================

%%K t10-apar-1 apar
\VUsaut{4.0mm}
\VUtitre{15.0pt}{60}{TABLICA Nr 10}
\VUsaut{3.0mm}
\VUpk{11.4pt}{40}{Morze. --- Port.}
\VUsaut{3.0mm}

%%K t10-01-1 p
1. --- Latem, podczas gdy jedni szukają spokoju i chłodu w górach, inni wolą udać się na
wybrzeże. Tak zrobiliśmy w zeszłym roku, bo bardzo lubimy \VUgras{Ocean}
\textsuperscript{(1)}.

%%K t10-02-1 p
2. --- Co do mnie, odczuwam wielką przyjemność, kontemplując ten ogromny bezmiar wody,
który rozciąga się aż po \VUgras{horyzont} \textsuperscript{(2)}, i widząc, jak odpływają
w daleką \VUgras{podróż} olbrzymie \VUgras{parowce} \textsuperscript{(3)}, na które
wsiadają podróżni jadący do Ameryki.

%%K t10-03-1 p
3. --- Dlatego mój ojciec, który zna moje upodobanie, zawsze wybiera na spędzenie
wakacji jakąś bardzo spokojną plażę, położoną blisko jednego z naszych wielkich
\VUgras{portów wojennych}.

%%K t10-03-2 p
Podczas naszego ostatniego pobytu \VUgras{eskadra} przybyła zakotwiczyć się za ogromnym
\VUgras{falochronem} \textsuperscript{(4)}, który zamyka i chroni \VUgras{port wojenny}
\textsuperscript{(5)}, i udało mi się zgodnie z życzeniem podziwiać rozmaite
\VUgras{okręty wojenne}, mały \VUgras{okręt podwodny} \textsuperscript{(10)}, który
nurkuje i znika pod wodą, a także ogromny \VUgras{pancernik} \textsuperscript{(9)},
który dotąd bał się prawie tylko ataków \VUgras{torpedowca} \textsuperscript{(8)}.

%%K t10-tit-2 sub
\VUsaut{3.0mm}
\VUpk{10.2pt}{40}{Małe okręty.}
\VUsaut{3.0mm}

%%K t10-04-1 p
4. --- Ten ostatni umiał już bardzo \VUgras{zręcznie} znaleźć słaby punkt grubego
metalowego pancerza, żeby wetknąć tam niebezpieczną \VUgras{torpedę}, której wybuch
sprawia zagładę okrętu, ale mimo to był mniej groźny niż okręt podwodny, bo
kontrtorpedowce łatwo go ścigają.

%%K t10-05-1 p
5. --- Mogłem też zobaczyć szybkie opancerzone \VUgras{krążowniki} \textsuperscript{(6)},
prawie tak samo nietykalne jak prawdziwy pancernik, kilka \VUgras{transportowców}
\textsuperscript{(12)} trzy albo cztery \VUgras{kanonierki} \textsuperscript{(7)} i
wreszcie jedną \VUgras{fregatę} \textsuperscript{(11)}. \VUgras{Widok tej floty, kiedy
manewruje, żeby podnieść kotwicę, jest najciekawszy.}

%%K t10-06-1 p
6. --- Jakaż różnica budowy między tymi wielkimi masami stali a eleganckimi
\VUgras{trójmasztowcami} czy wdzięcznymi \VUgras{żaglowcami} \textsuperscript{(13)},
używanymi jeszcze dziś we flocie handlowej, które widziałem wpływające do portu pod
pełnymi żaglami, kiedy spacerowałem po \VUgras{nabrzeżu} \textsuperscript{(14)}. Owszem,
traciły one wiele ze swoich zalet, \VUgras{kiedy} \VUgras{wiatr} był przeciwny. Musiały
zwinąć wszystkie żagle, żeby wrócić do basenu, a \VUgras{holownik}
\textsuperscript{(16)} był potrzebny, żeby po nie przypłynąć i doprowadzić je, niczym
zwykłe \VUgras{barki towarowe} \textsuperscript{(17)}, do nabrzeża.

%%K t10-tit-3 sub
\VUsaut{3.0mm}
\VUpk{10.2pt}{40}{Port handlowy.}
\VUsaut{3.0mm}

%%K t10-07-1 p
7. --- To, co jest ciekawe w porcie, o którym mówię, to że mieści on nie tylko port
wojenny, sztucznie urządzony olbrzymimi robotami, lecz także cudowny \VUgras{port
handlowy}, położony w \VUgras{ujściu} wielkiej rzeki.

%%K t10-08-1 p
8. --- Język ziemi, który oddziela oba baseny, jest ufortyfikowany i broniony szeregiem
\VUgras{baterii} i \VUgras{bastionów}, które wysuwają się w morze. Trzeba osobnego
pozwolenia, żeby spacerować po \VUgras{wałach} \textsuperscript{(15)}. Łatwo je
uzyskałem, przez mojego wuja, który jest inżynierem \VUgras{stoczni}
\textsuperscript{(18)}, i poszedłem obejrzeć okręty, które buduje się pod rozległymi
halami.

%%K t10-09-1 p
9. --- Byłem nawet przy wodowaniu pancernika i pamiętam wzruszającą chwilę, którą
poczuł cały tłum, kiedy ogromna masa, zostawiona samej sobie, zaczęła się ześlizgiwać po
swoim kołyskowym rusztowaniu i przyszła unosić się na wodzie rzeki.

%%K t10-10-1 p
10. --- Żeby dojść do stoczni, przechodzi się przez \VUgras{plac ćwiczeń}
\textsuperscript{(19)}, gdzie żołnierze garnizonu co dzień przychodzą się ćwiczyć,
przechodzi się po żelaznym \VUgras{mostku} \textsuperscript{(20)}, który łączy plac
defilad z \VUgras{arsenałem} \textsuperscript{(21)}.

%%K t10-tit-4 sub
\VUsaut{3.0mm}
\VUpk{10.2pt}{40}{Arsenał i doki.}
\VUsaut{3.0mm}

%%K t10-11-1 p
11. --- Z moim wujem zwiedziłem tę wielką budowlę pełną \VUgras{dział}
\textsuperscript{(22)}, \VUgras{kul} \textsuperscript{(23)} i \VUgras{granatów}.
Widziałem też \VUgras{fabrykę metalu} \textsuperscript{(24)}, w której wyrabia się
\VUgras{kotły} \textsuperscript{(25)} parowców.

%%K t10-12-1 p
12. --- Zupełnie blisko są \VUgras{doki} \textsuperscript{(26)} --- doki naprawcze --- i
bardzo osobliwe \VUgras{doki pływające} \textsuperscript{(27)}, w których mogłem
zobaczyć z bardzo bliska, na naprawianym okręcie, \VUgras{śrubę} \textsuperscript{(28)},
\VUgras{stępkę} \textsuperscript{(29)} i \VUgras{ster} \textsuperscript{(30)} statku.

%%K t10-13-1 p
13. --- Port handlowy jest zupełnie tak samo godny studiowania jak port wojenny i
spędziłem wiele godzin, gapiąc się na nabrzeżach, gdzie cumują \VUgras{statki kołowe}
\textsuperscript{(31)}, które płyną w górę rzeki, i patrząc, jak działają \VUgras{żurawie
masztowe} \textsuperscript{(32)}, które podnoszą, niczym piórka, ogromne ładunki.

%%K t10-tit-5 sub
\VUsaut{4.0mm}
\VUpk{10.2pt}{40}{Pilot.}
\VUsaut{3.0mm}

%%K t10-14-1 p
14. --- Podziwiałem \VUgras{transatlantyki} \textsuperscript{(34)} wychodzące z redy, z
\VUgras{flagami} \textsuperscript{(35)} na pełnym wietrze, prowadzone przez zręcznego
\VUgras{pilota} \textsuperscript{(36)}, który cudownie umie trzymać się \VUgras{kanaliku
żeglownego} znaczonego \VUgras{bojami} \textsuperscript{(40)} i \VUgras{pławami},
omijając \VUgras{galary} \textsuperscript{(41)}, które mozolnie dają się kierować swoim
jedynym czworokątnym żaglem. Rozróżniłem na \VUgras{przodzie statku}
\textsuperscript{(37)} --- dziobie --- \VUgras{kabiny} \textsuperscript{(39)} drugiej
klasy, a na \VUgras{tyle statku} \textsuperscript{(38)} --- rufie --- salony dla
pasażerów pierwszej klasy.

%%K t10-15-1 p
15. --- Kilka razy byłem też przy ładowaniu \VUgras{barki towarowej}
\textsuperscript{(42)}, w której właśnie gromadzono towary z \VUgras{magazynu}
\textsuperscript{(43)} i ze \VUgras{stacji morskiej} \textsuperscript{(44)}, przywiezione
przez liczne pociągi \VUgras{linii nabrzeżnej} \textsuperscript{(45)}.

%%K t10-tit-6 sub
\VUsaut{3.0mm}
\VUpk{10.2pt}{40}{Wioślarstwo dla przyjemności.}
\VUsaut{3.0mm}

%%K t10-16-1 p
16. --- Często pływałem łodzią w \VUgras{doku postojowym} \textsuperscript{(53)}, do
którego przychodzą się chronić \VUgras{łodzie} \textsuperscript{(46)}, i było dla mnie
radością władać \VUgras{wiosłem} \textsuperscript{(47)}. Wszedłem na \VUgras{pokład}
\textsuperscript{(48)} \VUgras{łodzi żaglowej} \textsuperscript{(49)}, wspiąłem się po
\VUgras{linach} \textsuperscript{(52)} na \VUgras{maszt} \textsuperscript{(51)} aż na
\VUgras{reje} \textsuperscript{(50)}, potem powtórzyłem swoją przechadzkę \VUgras{jolą}
\textsuperscript{(54)} między ciężkimi \VUgras{łodziami rybackimi}
\textsuperscript{(55)}, gdzie schną \VUgras{sieci} \textsuperscript{(56)}.

%%K t10-17-1 p
17. --- Na \VUgras{nabrzeżu} \textsuperscript{(58)} defilowały przede mną najrozmaitsze
\VUgras{towary} \textsuperscript{(59)}, zawarte w \VUgras{skrzyniach}
\textsuperscript{(60)} albo w \VUgras{belach} \textsuperscript{(61)}. Tworzyły
\VUgras{ładunek} \textsuperscript{(63)} \VUgras{barkasów}, a potężne \VUgras{dźwigi}
\textsuperscript{(62)} zdejmowały je i składały na \VUgras{wozach ciężarowych}
\textsuperscript{(64)}, podczas gdy \VUgras{woźnice} zgłaszali je i płacili cło w
\VUgras{komorze celnej} \textsuperscript{(65)}.

%%K t10-18-1 p
18. --- Wreszcie, kiedy przybyłem do \VUgras{mostu obrotowego} \textsuperscript{(66)}
albo do \VUgras{śluzy} \textsuperscript{(69)}, przechodząc przed \VUgras{pogłębiarką}
\textsuperscript{(68)}, która czyści \VUgras{kanał} \textsuperscript{(67)},
wysiadłem na nabrzeże blisko \VUgras{semafora} \textsuperscript{(70)}.

%%K t10-19-1 p
19. --- Po drugiej stronie tego nabrzeża przybijają właśnie \VUgras{szalupy}
\textsuperscript{(71)}, które przywożą na ląd oficerów eskadry. Godny podziwu to widok
patrzeć, jak marynarze rytmicznie podnoszą swoje wiosła, podczas gdy łódź niesiona przez
\VUgras{falę} \textsuperscript{(72)} łatwo dobija do schodków przystani. W tym miejscu
można lepiej obserwować \VUgras{pływ} i ruch \VUgras{przypływu} i \VUgras{odpływu} na
\VUgras{plaży} \textsuperscript{(73)}.

%%K t10-tit-7 sub
\VUsaut{3.0mm}
\VUpk{10.2pt}{40}{Kasyno oficerskie.}
\VUsaut{3.0mm}

%%K t10-20-1 p
20. --- Żeby wrócić do domu, skorzystałem z \VUgras{tramwaju elektrycznego}
\textsuperscript{(74)}, którego \VUgras{przystanek} \textsuperscript{(75)} jest blisko
\VUgras{słupa} postawionego przed klubem oficerów.

%%K t10-20-2 p
Z \VUgras{balkonu} \textsuperscript{(76)} klubu, ozdobionego \VUgras{kotwicami}
\textsuperscript{(77)}, działami i kulami, często widziałem wspaniałe zgromadzenie
oficerów marynarki: \VUgras{poruczników} \textsuperscript{(78)} ze szpadą,
\VUgras{kapitanów artylerii} \textsuperscript{(79)} i \VUgras{piechoty}
\textsuperscript{(80)} z \VUgras{szablą}. Za nimi stał \VUgras{marynarz}
\textsuperscript{(81)}, który przynosił im \VUgras{lunetę}, kiedy chcieli rozróżnić to,
co dzieje się daleko.

%%K t10-21-1 p
21. --- Widziałem też młodych \VUgras{podporuczników} \textsuperscript{(82)}
otrzymujących instrukcje swojego \VUgras{komendanta} \textsuperscript{(83)} albo
\VUgras{admirała} \textsuperscript{(84)}, podczas gdy \VUgras{bosman}
\textsuperscript{(85)} czekał, żeby zanieść je na okręt.

%%K t10-22-1 p
22. --- Z tego balkonu widok jest wspaniały: panuje się nad całą plażą z jej
\VUgras{namiotami} \textsuperscript{(86)} i \VUgras{kabinami} \textsuperscript{(87)}, i
widzi się, o godzinie kąpieli, \VUgras{kąpiących się} \textsuperscript{(88)}, którzy
radośnie swawolą przed \VUgras{kasynem} \textsuperscript{(89)}.

%%K t10-23-1 p
23. --- Na krańcu plaży wybrzeże tworzy mały skalisty \VUgras{półwysep}
\textsuperscript{(91)}, gdzie rozbija się \VUgras{fala przybojowa} \textsuperscript{(92)}
i na którym zbudowano \VUgras{latarnię morską} \textsuperscript{(93)}, nocą wskazującą
drogę żeglującym. Ten półwysep łączy się z lądem tylko bardzo wąskim \VUgras{przesmykiem}
\textsuperscript{(90)}.

%%K t10-tit-8 sub
\VUsaut{3.0mm}
\VUpk{10.2pt}{40}{Ogólny widok wybrzeży.}
\VUsaut{3.0mm}

%%K t10-24-1 p
24. --- \VUgras{Klif} \textsuperscript{(94)} rozciąga się, rozdarty przez morze, które
kapryśnie rysuje w nim szereg \VUgras{zatok} \textsuperscript{(95)} i
\VUgras{przylądków} (cypli), czyniąc go najbardziej malowniczym.

%%K t10-24-2 p
Na \VUgras{płaskowyżu} \textsuperscript{(96)}, który panuje nad \VUgras{(morską)
cieśniną} \textsuperscript{(98)}, jest bardzo ważny \VUgras{fort} \textsuperscript{(97)}
i kilka «willi».

%%K t10-25-1 p
25. --- Ta cieśnina oddziela od kontynentu bardzo niebezpieczną \VUgras{wyspę}
\textsuperscript{(99)}, otoczoną \VUgras{rafami} \textsuperscript{(100)} i
\VUgras{przybojem}, gdzie łodzie, a nawet wielkie okręty czasem osiadają na mieliźnie.
Mimo szybkości ratunku i prędkiego przybycia \VUgras{łodzi ratunkowych}, te
\VUgras{katastrofy morskie} są przerażające, a podczas \VUgras{burz} równonocnych
kilka statków zginęło z ludźmi i towarami.

%%K t10-25-2 p
Mimo tego sąsiedztwa \VUgras{port} jest bardzo uczęszczany przez marynarzy.
\VUsaut{6.0mm}
\VUfilet{22mm}
